\documentclass[11pt]{article}

\usepackage[margin=1in]{geometry}
\usepackage{amsmath,amssymb,bm}
\usepackage{natbib}
\usepackage{microtype}
\usepackage[T1]{fontenc}
\usepackage{lmodern}


% Macros used in the supplementary material
\newcommand{\vy}{\boldsymbol{y}}
\newcommand{\veta}{\boldsymbol{\eta}}

\begin{document}

\begin{center}
{\Large
Supplementary material to \\ ``The Lasso Distribution: Properties, Sampling Methods, and Applications in Bayesian Lasso Regression''}

\bigskip 

 
by Mohammad Javad Davoudabadi, Jonathon Tidswell, Samuel Muller,  Garth Tarr \\ and John T. Ormerod
 

\bigskip 
\date{\today}
\end{center}

\appendix

\section{Derivation of the full conditional for each $\beta_j$}\label{Suppl:derivation_fullConditional_Dist}
Consider $p({\bf y}| \boldsymbol{\beta}, {\veta})$, $p(\boldsymbol{\beta}| {\veta})$, $p({\veta})$, $p(\sigma^2)$, and $p(\lambda^2)$ are the likelihood function, and the priors of $\boldsymbol{\beta}$, ${\veta}$, $\sigma^2$, and $\lambda^2$, respectively. The log of the full joint distribution is proportional to
\begin{align*}
    \log p(\vy, \boldsymbol{\beta}, \sigma^2, \lambda^2,
{\veta}) & = \log \big(p({\bf y}| \boldsymbol{\beta,{\veta}})p(\boldsymbol{\beta}|{\veta})p({\veta})p(\sigma^2)p(\lambda^2) \big)\\ &
\propto - \frac{1}{2\sigma^2} \Big[ {\bf y}^T{\bf y} - 2{\bf y}^T \mathbf{X}^T \boldsymbol{\beta} + \boldsymbol{\beta}^T \mathbf{X}^T \mathbf{X} \boldsymbol{\beta} + \lambda^2 \boldsymbol{\beta}^T{\bf A}\boldsymbol{\beta} \Big]  \\ 
& \quad\quad{}- \frac{\sum_{i=1}^p \eta_i}{2} - 2\sum_{i=1}^p \eta_i - \sum_{i=1}^p\frac{1}{2\eta_i} + \mbox{constant}; 
\end{align*}
where ${\veta}= (\eta_1,\ldots,\eta_p)$ and ${\bf A}= \text{diag}({\veta}^{-1})$. After integrating out ${\veta}$ from the above proportion, the marginal distribution of $\beta_j$ is as follows
\begin{align*}
    \log p(\beta_j| \text{rest}) \propto - \frac{\parallel  \mathbf{X}_j\parallel^2_2}{2\sigma^2}\beta_j^2 + \frac{\mathbf{X}_j^T({\bf y} - \mathbf{X}_{-j}\boldsymbol{\beta}_{-j})}{\sigma^2} \beta_j - \frac{\lambda^2}{\sigma} |\beta_j|;
\end{align*}
which is a univariate Lasso distribution with parameters $a = \frac{\parallel  \mathbf{X}_j\parallel^2_2}{\sigma^2}$, $b = \frac{\mathbf{X}_j^T({\bf y} - \mathbf{X}_{-j}\boldsymbol{\beta}_{-j})}{\sigma^2}$, and $c = \frac{\lambda^2}{\sigma}$.



\section{Derivation of the distributional properties of the Lasso distribution}\label{Suppl:deriv_LassoDist_properties}
If $X \sim \mbox{Lasso}(a,b,c)$ with then it has density given by
$$
p(X,a,b,c) = Z^{-1}\exp\left( -\tfrac{1}{2}aX^2 + bX - c|X| \right)
$$
where $Z$ is the normalizing constant. Then
$$
\begin{array}{rl}
Z(a,b,c)
& = \int_{-\infty}^\infty \exp\left[ -\tfrac{1}{2}aX^2 + bX - c|X| \right] dX
\\ [2ex]
& 
= \int_0^\infty    \exp\left[ -\tfrac{1}{2}aX^2 + (b - c)X \right] dx
+ \int_{-\infty}^0 \exp\left[ -\tfrac{1}{2}aX^2 + (b + c)X \right] dX
\\ [2ex]
&  
= \int_0^\infty \exp\left[ - \frac{(X - \mu_1)^2}{2\sigma^2} + \frac{\mu_1^2}{2\sigma^2} \right] dx
+ \int_{-\infty}^0 \exp\left[ - \frac{(X - \mu_2)^2}{2\sigma^2} + \frac{\mu_2^2}{2\sigma^2} \right] dX
\\ [2ex]
&  
= \sqrt{2\pi\sigma^2}
\left[  \exp\left\{  \frac{\mu_1^2}{2\sigma^2} \right\} \int_0^\infty \phi(x;\mu_1,\sigma^2) dx
+       \exp\left\{  \frac{\mu_2^2}{2\sigma^2} \right\} \int_{-\infty}^0 \phi(x;\mu_2,\sigma^2) dx
\right] 
\\ [2ex]
&  
= \sqrt{2\pi\sigma^2}
\left[  \exp\left\{  \frac{\mu_1^2}{2\sigma^2} \right\} \left\{ 1 - \Phi(-\mu_1/\sigma) \right\} 
+       \exp\left\{  \frac{\mu_2^2}{2\sigma^2} \right\} \left\{\Phi(-\mu_2/\sigma) \right\} 
\right] 
\\ [2ex]
&  
= \sqrt{2\pi\sigma^2}
\left[  \exp\left(  \frac{\mu_1^2}{2\sigma^2} \right) \Phi\left(\frac{\mu_1}{\sigma} \right) 
+       \exp\left(  \frac{\mu_2^2}{2\sigma^2} \right) \Phi\left( \frac{-\mu_2}{\sigma} \right)  
\right] 
\\ [2ex]
& 
= 
  \sigma \left[ \frac{\Phi(\mu_1/\sigma)}{\phi(\mu_1/\sigma)}
+ \frac{\Phi(-\mu_2/\sigma)}{\phi(-\mu_2/\sigma)}  \right] 
\end{array} 
$$
where $\mu_1 = (b-c)/a$, $\mu_2 = (b + c)/a$ and $\sigma^2 = 1/a$.

The MGF of the Lasso distribution is
\begin{align*}
    M(t) = \mathbb{E} (\exp (tX)) & = Z(a,b,c)^{-1} \int \exp (tX) \exp \left(\frac{-a}{2}X^2 + bX -c|X| \right) dX \\ &
    = Z(a,b,c)^{-1} \int \exp \left(\frac{-a}{2}X^2 + (b+t) X -c|X|\right) dX \\ &
    = \frac{Z(a,b+t,c)}{Z(a,b,c)}Z(a,b+t,c)^{-1} \int \exp \left(\frac{-a}{2}X^2 + (b+t) X -c|X|\right) dX  \\ &
    = \frac{Z(a,b+t,c)}{Z(a,b,c)}.
\end{align*}

The moments of the Lasso distribution are:
$$
\begin{array}{rl}
\mathbb{E}(X^r)
& = Z^{-1} \int_{-\infty}^\infty X^r \exp\left[ -\tfrac{1}{2}aX^2 + bX - c|X| \right] dX
\\ [2ex]
&  
= Z^{-1}  \int_0^\infty   X^r \exp\left[ -\tfrac{1}{2}aX^2 + (b - c)X \right] dX
+ \int_{-\infty}^0 X^r \exp\left[ -\tfrac{1}{2}aX^2 + (b + c)X \right] dX
\\ [2ex]
&  
= Z^{-1}  \sqrt{2\pi\sigma^2}
 \exp\left(  \frac{\mu_1^2}{2\sigma^2} \right) \int_0^\infty x^r \phi(x;\mu_1,\sigma^2) dx
 +   \sqrt{2\pi\sigma^2}   \exp\left(  \frac{\mu_2^2}{2\sigma^2} \right) \int_{-\infty}^0 x^r \phi(x;\mu_2,\sigma^2) dx

\\ [2ex]
& 
= \frac{\sigma}{Z} \left[  
 \frac{\Phi(\mu_1/\sigma)}{\phi(\mu_1/\sigma)} \frac{\int_0^\infty X^r \phi(x;\mu_1,\sigma^2) dx}{\Phi(\mu_1/\sigma)}
+  \frac{\Phi(-\mu_2/\sigma)}{\phi(-\mu_2/\sigma)}  \frac{\int_{-\infty}^0 X^r \phi(x;\mu_2,\sigma^2) dx}{\Phi(\mu_2/\sigma)}
\right] 

\\ [4ex]
&  
= \frac{\sigma}{Z} \left[  
 \frac{\Phi(\mu_1/\sigma)}{\phi(\mu_1/\sigma)} 
 \mathbb{E}( A^r )
+ \frac{\Phi(-\mu_2/\sigma)}{\phi(-\mu_2/\sigma)}  \mathbb{E}( B^r )
\right] 

\\ [4ex]
&  
= \frac{\sigma}{\sigma \left[ \frac{\Phi(\mu_1/\sigma)}{\phi(\mu_1/\sigma)}
+ \frac{\Phi(-\mu_2/\sigma)}{\phi(-\mu_2/\sigma)}  \right]} \left[  
 \frac{\Phi(\mu_1/\sigma)}{\phi(\mu_1/\sigma)} 
 \mathbb{E}( A^r )
+  \frac{\Phi(-\mu_2/\sigma)}{\phi(-\mu_2/\sigma)}  \mathbb{E}( B^r )
\right] 
\\ [4ex]
&  
= \left[ \frac{\frac{\Phi(\mu_1/\sigma)}{\phi(\mu_1/\sigma)}}{\left[ \frac{\Phi(\mu_1/\sigma)}{\phi(\mu_1/\sigma)}
+ \frac{\Phi(-\mu_2/\sigma)}{\phi(-\mu_2/\sigma)}  \right]}    
 \mathbb{E}( A^r )
+  \frac{\frac{\Phi(-\mu_2/\sigma)}{\phi(-\mu_2/\sigma)}}{\left[ \frac{\Phi(\mu_1/\sigma)}{\phi(\mu_1/\sigma)}
+ \frac{\Phi(-\mu_2/\sigma)}{\phi(-\mu_2/\sigma)}  \right]}   \mathbb{E}( B^r )
\right] 
\\ [4ex]
&  
=
 (1 - w) \mathbb{E}(A^r) +  w \mathbb{E}(B^r)

\end{array} 
$$
where $A\sim TN_+(\mu_1,\sigma^2)$, $B\sim TN_-(\mu_2,\sigma^2)$ and $TN_+$ and
 $TN_-$ denote the positively and negatively truncated
 normal distributions respectively.

If $x\le 0$ the CDF is given by
$$
\begin{array}{rl} 
P(X < x) 
&  = Z^{-1} \int_{-\infty}^x \exp\left[ -\tfrac{1}{2}at^2 + (b + c)t \right] dt 
\\ [2ex]
& = Z^{-1} \sqrt{2\pi\sigma^2} \exp\left( \frac{\mu_2^2}{2\sigma^2} \right) \int_{-\infty}^x \phi(t;\mu_2,\sigma^2) dt
\\ [2ex]
& = Z^{-1} \sqrt{2\pi\sigma^2} \exp\left( \frac{\mu_2^2}{2\sigma^2} \right) \Phi\left( \frac{x - \mu_2}{\sigma} \right)
\\ [2ex]
& =  \frac{\sigma}{Z}\frac{ \Phi\left( \frac{x - \mu_2}{\sigma} \right)}{\phi(-\mu_2/\sigma)}
\\ [2ex]
& = w\frac{\Phi\left( 
\frac{x - \mu_2}{\sigma}\right)}{\Phi(-\mu_2/\sigma)} 

\end{array} 
$$
and if $x>0$ we have
$$
\begin{array}{rl}
P(X < x)
&  = Z^{-1}\int_{-\infty}^x \exp\left( -\tfrac{1}{2}at^2 + bt - c|t| \right) dt
\\ [2ex]
& 
= Z^{-1} \sqrt{2\pi\sigma^2}
\left[  \exp\left(  \frac{\mu_1^2}{2\sigma^2} \right) \int_0^x \phi(t;\mu_1,\sigma^2) dt
+      \exp\left(  \frac{\mu_2^2}{2\sigma^2} \right) \Phi\left( \frac{-\mu_2}{\sigma} \right)  
\right] 
\\ [2ex]
&  = \frac{\sigma}{Z} 
\left[  \frac{ 
\Phi\left( \frac{x - \mu_1}{\sigma} \right) - \Phi\left( \frac{ - \mu_1}{\sigma} \right)}{\phi(\mu_1/\sigma)}
+       \frac{\Phi\left( -\mu_2/\sigma \right)}{\phi(-\mu_2/\sigma)}  
\right] 
\\ [2ex]
&  = \displaystyle 1 - (1 - w)\frac{\Phi(\frac{\mu_1 - x}{\sigma})}{\Phi(\mu_1/\sigma)}.

\end{array} 
$$
For the inverse CDF, we again have two cases. Let $u = P(X<x)$. When 
$$
u \le \frac{\sigma}{Z}  \frac{ \Phi\left(  - \mu_2/\sigma \right)}{\phi(-\mu_2/\sigma)} = w
$$ 
we solve
$$
u = \frac{ \sigma \Phi\left( \frac{x - \mu_2}{\sigma} \right)}{Z \phi(-\mu_2/\sigma)}
$$
for $x$ to obtain
$$
x = \mu_2 + \sigma \,\Phi^{-1}\left[ (Z/\sigma) \phi(-\mu_2/\sigma) u \right] = \mu_2 + \sigma \,\Phi^{-1}\left[ \Phi(-\mu_2/\sigma) u /w\right].
$$ 
When 
$$
u > \frac{\sigma}{Z}  \frac{ \Phi\left(  - \mu_2/\sigma \right)}{\phi(-\mu_2/\sigma)} = w
$$ 
we need to solve
$$
u = \displaystyle 1 - (1 - w)\frac{\Phi(\frac{\mu_1 - x}{\sigma})}{\Phi(\mu_1/\sigma)}
$$
for $x$ to obtain
$$
x = \mu_1 - \sigma\, \Phi^{-1}\left[ \Phi(\mu_1/\sigma) (1-u)/(1-w) \right]. 
$$

The Lasso distribution belongs to the exponential family distributions based on the following in which ${\boldsymbol\theta} = (a, b, c)$ is the set of parameters
\begin{align*}
p(X|{\boldsymbol\theta})
    & = h(X)\exp\left[{\boldsymbol\nu}({\boldsymbol\theta}) \textbf{T}(X)-\mathbb{A}({\boldsymbol\theta})\right]\\ 
%    & =Z^{-1}\exp(-\frac{1}{2}ax^2+bx-c|x|)\\&
%    =\sigma^* \left[\frac{\Phi(\mu_1/\sigma^*)}{\phi(\mu_1/\sigma^*)} + \frac{\Phi(-\mu_2/\sigma^*)}{\phi(-\mu_2/\sigma^*)} \right]^{-1} \exp(-\frac{1}{2}ax^2+bx-c|x|)\\&
    & = \exp\left[ -\tfrac{1}{2}aX^2+bX-c|X| - \left\{ \log\left(\frac{\Phi(\mu_1/\sigma)}{\phi(\mu_1/\sigma)} + \frac{\Phi(-\mu_2/\sigma)}{\phi(-\mu_2/\sigma)} \right) + \log (\sigma)  \right\} \right], 
\end{align*}
where
$$
\begin{array}{c}
{\boldsymbol\nu}({\boldsymbol\theta}) 
    = (-\frac{1}{2}a,b,-c), 
    \quad
h(X) = 1, 
    \quad 
\textbf{T}(X) = (X^2,X,|X|), \quad \text{and}
\\
\displaystyle \mathbb{A}({\boldsymbol\theta}) 
    = \log\left[ \frac{\Phi(\mu_1/\sigma)}{\phi(\mu_1/\sigma)} + \frac{\Phi(-\mu_2/\sigma)}{\phi(-\mu_2/\sigma)} \right] + \log (\sigma).
\end{array}
$$


\paragraph{Mode of the Lasso distribution.}
Since the normalizing constant does not depend on $x$, the mode maximizes the log-kernel
\[
\ell(x)= -\tfrac12 a x^2 + b x - c|x|,\qquad a>0.
\]
The function $\ell(x)$ is concave and piecewise quadratic. For $x\ge 0$ we have
\[
\ell(x)= -\tfrac12 a x^2 + (b-c)x,
\]
whose maximizer is $x^*=(b-c)/a$, truncated to the feasible set $x\ge 0$, hence
\[
x^*_+ = \max\{(b-c)/a,\,0\}.
\]
For $x\le 0$ we have
\[
\ell(x)= -\tfrac12 a x^2 + (b+c)x,
\]
whose maximizer is $x^*=(b+c)/a$, truncated to $x\le 0$, hence
\[
x^*_- = \min\{(b+c)/a,\,0\}.
\]
Comparing the two cases (equivalently, applying the soft-thresholding operator) yields the unique global maximizer
\[
\operatorname{mode}(X)= \frac{\operatorname{sign}(b)}{a}\,\max\{|b|-c,\,0\}.
\]
In particular, if $|b|\le c$ then $0$ is the mode, while if $|b|>c$ the mode is shifted from $b/a$ toward $0$ by $c/a$.

\section{Numerical accuracy of Mills ratio approximation}
We verified the numerical accuracy of the implemented Mills ratio function by comparing it against a 256-bit arbitrary-precision reference computed using the \texttt{Rmpfr} package \citep{Rmpfr}. Over the range $d \in [0,2000]$, the maximum relative error was $6.31 \times 10^{-14}$ and the median relative error was $1.25 \times 10^{-14}$, corresponding to at least 13 correct significant digits across all tested values (see Table~\ref{tab:mills_accuracy}). 

For comparison, the naive evaluation $ \Phi(-d)/\phi(-d)$ using standard double-precision routines (\texttt{pnorm}/\texttt{dnorm}) becomes numerically unstable for $d \gtrsim 37.5$, where underflow leads to $\texttt{NaN}$ or infinite values. In contrast, the proposed rational approximation remains stable and highly accurate throughout the entire tested range.

\begin{table}[ht]
\centering
\begin{tabular}{lcc}
\hline
Metric & Value \\ 
\hline
Maximum relative error & $6.31 \times 10^{-14}$ \\
Median relative error  & $1.25 \times 10^{-14}$ \\
Minimum significant digits & $13.20$ \\
Median significant digits  & $13.90$ \\
\hline
Naive evaluation failure threshold & $d \approx 37.5$ \\
\hline
\end{tabular}
\caption{Accuracy of the implemented Mills ratio approximation compared to a 256-bit arbitrary-precision reference. Relative errors are computed over selected points in $d \in [0,2000]$. Significant digits are estimated as $-\log_{10}(\text{relative error})$.}
\label{tab:mills_accuracy}
\end{table}


\bibliographystyle{plainnat}
\bibliography{MJ}


\end{document}
